\section{Introduction}
\label{sec:intro}

\IEEEPARstart{F}{loods} disrupt road access exactly when isolated communities need supplies most. Small uncrewed aerial vehicles (UAVs) can reach such locations quickly and cheaply, and the disaster-management literature repeatedly identifies mapping, situational awareness and the delivery of small, urgent payloads as their most promising roles \cite{erdelj2017help,mohddaud2022applications,khan2022emerging,rejeb2021humanitarian}. Flood extents are now routinely produced from satellite and aerial data \cite{twele2016sentinel,bonafilia2020sen1floods11,rahnemoonfar2021floodnet} and are widely published as annotated maps. Turning such a map into a flight plan, however, is still largely a manual task: an operator must decide where packages can be dropped, in which order to visit those points, how to route around water, and how to encode the result in the mission format of an autopilot.

Each of these steps has a mature research literature, but the steps are usually studied in isolation. Flood segmentation methods stop at a mask \cite{gebrehiwot2019deep,simantiris2024unsupervised}; flood inundation models require digital elevation models and hydraulic forcing \cite{bates2000simple,ghimire2013formulation,teng2017flood}; relief-logistics models assume that demand locations are already known nodes in a graph \cite{otto2018optimization,rabta2018drone,dong2026path}; and path planners assume that an obstacle map is given \cite{aggarwal2020path,koenig2005fast}. A recent review of adaptive mission planning in emergencies observes that methods coupling hazard evolution, observation and mission planning remain fragmented across application domains \cite{gioia2026perspective}. Weather adds a further dynamic constraint, both on whether a drone can fly at all \cite{gao2021weather} and on how a flood may evolve.

\subsection{Problem Statement and Research Question}
Given a single flood-annotated map raster, a geographic reference and the current weather, the task is to produce, without training data or terrain models, a UAV mission that (i) visits one accessible delivery point for every flooded region, (ii) avoids current and near-term flooded areas where possible, (iii) releases a payload at each point, and (iv) can be loaded by a MAVLink-based ground station. When several UAVs are available, the same task is to distribute the delivery points among vehicles operating from separate launch bases, so that every vehicle's sorties fit its battery and payload and the vehicles' routes do not cross. The research question of this paper is therefore: \emph{to what extent can a lightweight, training-free pipeline convert an annotated flood map and current weather into an executable and verifiably correct UAV relief mission, and which limitations prevent such a mission from being geographically and operationally trustworthy?}

\subsection{Two Planning Approaches}
\adapt{} offers two ways of serving the same set of delivery points (Fig.~\ref{fig:modes}).
\begin{itemize}
\item \emph{Approach A: one drone.} A single UAV takes off from one home point, visits every delivery point in one tour and returns. This is the single-UAV mode (Section~\ref{sec:method}).
\item \emph{Approach B: a fleet of drones.} Up to $K$ UAVs operate from separate launch bases. The planner places the fewest bases that can reach the delivery points within battery range, then gives every delivery point to the nearest base that can serve it. Each drone therefore serves the cluster of flood regions around its own base, flying as many battery-limited sorties as it needs. This is the multi-base mode (Section~\ref{sec:multi}).
\end{itemize}
With $K=1$ the same planner produces a one-drone mission, which allows a like-for-like comparison of the two approaches (Section~\ref{sec:multires}). Approach B resembles what is often called a drone swarm, but in \adapt{} the drones are coordinated only before flight. The planner divides the work and checks that the routes do not cross; the drones do not communicate or adapt to each other while flying.

\subsection{Objectives}
The objectives are to (O1) describe and formalise every stage of an existing open implementation, \adapt{}, exactly as implemented; (O2) verify the correctness of its geospatial and mission-generation layers independently of the implementation; (O3) characterise its behaviour on real published flood maps, including ablation, sensitivity, optimality and baseline comparisons, and compare single- and multi-UAV planning on identical missions; and (O4) state explicitly what the evidence does and does not establish.

\subsection{Contributions}
The contributions are deliberately limited to what the implementation and its evaluation support:
\begin{enumerate}
\item \textbf{An integrated image-to-mission pipeline.} \adapt{} chains colour-based flood extraction, a weather-parameterised obstacle-inflation rule, boundary drop-point selection, nearest-neighbour ordering, per-leg D*~Lite planning and pixel-to-geodetic conversion into a QGroundControl WPL~110 mission with explicit payload-release commands (Sections~\ref{sec:arch}--\ref{sec:method}). The individual algorithms are established; the contribution is their integration and the resulting executable artefact.
\item \textbf{A single-drone and a drone-fleet approach.} Besides the single-UAV pipeline, a multi-base mode works as follows (Section~\ref{sec:multi}): it selects launch bases, gives each drone the cluster of delivery points nearest its base under a battery and payload model, plans one mission per drone and rejects plans whose routes cross. A controlled comparison on identical missions isolates the effect of the number of vehicles. The mode is a pre-flight planner with a nearest-feasible assignment, not a cooperative swarm.
\item \textbf{A precise formulation with stated validity domains.} Each stage is specified as implemented, including the failure behaviour of the planner and the conditions under which the geodetic conversion refuses to produce a coordinate (Section~\ref{sec:method}).
\item \textbf{A verification methodology for the geospatial and mission layers.} An independent WGS84 oracle, a specification-based mission validator that is itself tested against corrupted files, metamorphic resize tests and mutation testing are used to establish software correctness separately from geographic validity (Sections~\ref{sec:verif} and~\ref{sec:results}).
\item \textbf{An empirical characterisation that reports unfavourable outcomes.} On three published flood maps we report end-to-end behaviour, an ablation of predicted-flood obstacles, threshold sensitivity, the optimality gap of the route ordering, a planner baseline and a single- versus multi-UAV comparison, including failure modes (Section~\ref{sec:results}).
\end{enumerate}
We do not claim real-time operation, geographic accuracy for the supplied maps, flood-prediction accuracy, or simulated or physical flight validation; Section~\ref{sec:limits} discusses why.

The remainder of the paper is organised as follows. Section~\ref{sec:related} reviews related work and derives the gap. Section~\ref{sec:arch} formalises the problem and the architecture, Section~\ref{sec:method} details the single-UAV methodology, and Section~\ref{sec:multi} describes the multi-UAV mode. Section~\ref{sec:results} presents the experimental setup and results. Section~\ref{sec:limits} discusses limitations, threats to validity and safety. Section~\ref{sec:compare} positions \adapt{} relative to prior work, Section~\ref{sec:future} outlines future work, and Section~\ref{sec:conclusion} concludes.
